\documentclass{article}

\usepackage{amsmath}
\usepackage{amssymb}
\usepackage{amsthm}
\usepackage{listings}

\newtheorem{lemma}{Lemma}
\newtheorem{definition}{Definition}

\lstset{basicstyle=\ttfamily}

%% \case is already defined in the paper (by lncs?)
\newcommand{\Case}[2]{\textbf{Case #1:} #2}
\newcommand{\basecase}[1]{\textbf{Base Case:} #1}
\newcommand{\inductionStep}{\textbf{Induction Step.}}
\newcommand{\wfdomn}{\mathbb{V}^{n}}
\newcommand{\wfdom}[1]{\mathbb{V}^{#1}}

\newcommand{\ith}[1]{#1[i]}
\newcommand{\lsb}[1]{#1[0]}
\newcommand{\tv}[1]{#1.v}
\newcommand{\tm}[1]{#1.m}

\newcommand{\ithv}[1]{\ith{\tv{#1}}}
\newcommand{\ithm}[1]{\ith{\tm{#1}}}

\newcommand{\wf}[1]{wellformed(#1)}
\newcommand{\ingamma}[2]{#1 \in \gamma(#2)}

\newcommand{\tadd}[2]{tnum\_add(#1, #2)}
\newcommand{\tunion}[2]{tnum\_union(#1, #2)}

\newcommand{\seti}[1]{set(#1, i)}

\newcommand{\sv}[2]{(\tv{#1} + \tv{#2})}
\newcommand{\sm}[2]{(\tm{#1} + \tm{#2})}
\newcommand{\svm}[1]{(\tv{#1} + \tm{#1})}
\newcommand{\orvm}[1]{(\tv{#1} \mathbin{|} \tm{#1})}

\begin{document}

\title{Tnum Arithmetic Proofs for Humans}
\author{Nandakumar Edamana}

\maketitle

\begin{abstract}
  This document provides a pen-and-paper-style explanation for
  our Rocq proofs related to tnum arithmetic.
\end{abstract}

\section{Convention}

This document uses upper case letters to denote tnums that have a
value-mask representation internally.  $\wfdomn$ is used to represent
the set of all possible well-formed (see
Definition~\ref{def:wellformed}) tnums in $n$-bit value-mask
representation.  For a tnum $P$, $\ithv{P}$ represents the $i$-th
value bit of $P$, and $\ithm{P}$ represents the $i$-th mask bit of $P$
(like in the Rocq proofs, this document uses lower numbers to
represent least significant bits, but 1-indexing is used instead of
0-indexing). For a bit vector $x$, $\seti{x}$ denotes the same bit
vector with its $i$-ith bit set to $1$.

\section{Background}

In this section, $i$ denotes a valid index (bit position) for the bit
vector it is used with. Details like bounds are omitted for brevity.

\begin{definition}
  \label{def:wellformed} (Well-formedness)
  $$\forall P \in \wfdomn, \ithm{P} = 1 \implies \ithv{P} = 0$$
\end{definition}

\begin{definition}
  \label{def:ingamma} (Membership)
  $$\forall p \in \{0, 1\}^{n}, P \in \wfdomn : \ingamma{p}{P} \iff (\ithm{P} = 0 \implies \ith{p} = \ithv{P})$$
\end{definition}

\begin{lemma}
  \label{lem:ingamma_value}
  $\forall P \in \wfdomn, \ingamma{\tv{P}}{P}$
\end{lemma}
\begin{proof}
  Follows from Definition~\ref{def:ingamma}.
\end{proof}

\begin{lemma}
  \label{lem:ingamma_value_bitor_mask}
  $\forall P \in \wfdomn, \ingamma{\orvm{P}}{P}$
\end{lemma}
\begin{proof}
  Follows from Definition~\ref{def:ingamma}.
\end{proof}

\begin{lemma}
  \label{lem:ingamma_set_at_mask}
  Consider any $P \in \wfdomn$. If $\ithm{P} = 1$, then $\ingamma{\seti{\tv{P}}}{P}$.
\end{lemma}
\begin{proof}
  Follows from Definition~\ref{def:ingamma}. (Note that
  $\seti{\tv{P}}$ is the same as $\tv{P}$, except at position $i$. At
  position $i$, the mask is set, so the value bit is free to be either
  $0$ or $1$ without losing membership.)
\end{proof}

\begin{lemma}
  \label{lem:tnum_add_v_m_is_or}
  $\forall P \in \wfdomn, \tv{P} + \tm{Q} = \tv{P} \mathbin{|} \tm{Q}$
\end{lemma}
\begin{proof}
  Due to well-formedness, $\ithv{P}$ and $\ithm{P}$ cannot be both $1$
  for any $i$. For such inputs, addition decomposes to bitwise OR.
\end{proof}

\section{Optimality of \lstinline{tnum_add}}

Informally speaking, an abstract arithmetic operation is optimal if it
has no more uncertainty than there is in the concrete result set.

In the case of tnum addition, this means if the abstract sum of two
tnums have $\mu$ at some position $i$, there should be at least one
pair of concrete sums that differ at $i$. We formally state and prove
this in the following lemma, showing \lstinline{tnum_add} from the
Linux kernel is optimal.

\begin{lemma}
  Consider any $P, Q \in \wfdomn$, and a position $i \le n$. If $\ithm{\tadd{P}{Q}} = 1$, then
  \[
  \exists \ingamma{p, p'}{P}, \exists \ingamma{q, q'}{Q} : \ith{(p + q)} \neq \ith{(p' + q')}.
  \]
\end{lemma}

\begin{proof}
  There are exactly three cases by which $\ithm{\tadd{P}{Q}}$ can be
  $1$. Either $\ithm{P} = 1$, $\ithm{Q} = 1$, or $\ithm{\chi} = 1$.
  We consider these three cases individually and show that there exist
  $\ingamma{p, p'}{P}$ and $\ingamma{q, q'}{Q}$ such that $\ith{(p +
    q)} \neq \ith{(p' + q')}$. (Note that either of $p \neq p'$ or $q
  \neq q'$ must hold, but both need not be true.)

  \Case{1}{$\ithm{P} = 1$}. Taking $p = \tv{P}$, $p' = \seti{\tv{P}}$,
  and $q = q' = \tv{Q}$ satisfies our goal. Their membership follows
  from Lemma~\ref{lem:ingamma_value} and
  Lemma~\ref{lem:ingamma_set_at_mask}. We know that $\ithv{P} = 0$
  since {$\ithm{P} = 1$}. Hence $p$ and $p'$ differ at (and only at)
  position $i$, and $\ith{(p + q)} \neq \ith{(p' + q')}$.

  \Case{2}{$\ithm{Q} = 1$}. Similar to the previous case, we take $q =
  \tv{Q}$, $q' = \seti{\tv{Q}}$, and $p = p' = \tv{P}$.

  \Case{3}{$\ithm{\chi} = 1$}. Recall that $\chi = (\sv{P}{Q} +
  \sm{P}{Q}) \oplus \sv{P}{Q}$. Since we know that $\ithm{\chi} = 1$,
  we have two possibilities here:

  \Case{3.1}{$\ith{(\sv{P}{Q} + \sm{P}{Q})} = 1, \ithv{P} = \ithm{P} = \ith{\sv{P}{Q}} = 0$}.

  After a basic rearranging and application of
  Lemma~\ref{lem:tnum_add_v_m_is_or}, $\ith{(\orvm{P} + \orvm{Q})} =
  1$. We know that $\ingamma{\orvm{P}}{P}$ and $\ingamma{\orvm{Q}}{Q}$
  (by Lemma~\ref{lem:ingamma_value_bitor_mask}). Hence we can pick
  $\ingamma{p, p'}{P}$ and $\ingamma{q, q'}{Q}$ in the following
  manner to satisfy $\ith{(p + q)} \neq \ith{(p' + q')}$:

  \begin{align*}
    p &= \tv{P} \\
    q &= \tv{Q} \\
    p' &= \orvm{P} \\
    q' &= \orvm{Q}
  \end{align*}

  \Case{3.2}{$\ith{\sv{P}{Q}} = 1, \ith{(\sv{P}{Q} + \sm{P}{Q})} =
    \ithv{P} = \ithm{P} = 0$}. Same assignments for $p$, $p'$, $q$ and
  $q'$ work in this case as well.
\end{proof}

As a side note, it is interesting to note that $\tv{P}$ is the
smallest concrete value and $\orvm{P}$ is the largest concrete value
represented by $P$, for any $P$.

\section{Optimality of \lstinline{tnum_union}}

The concrete union of two tnums $P$ and $Q$ is $\gamma(P) \cup
\gamma(Q)$. This can be rewritten like in the following:

\begin{definition}
  \[
  \text{concrete\_union}(P, Q) = \{ x \mathbin{|} \ingamma{x}{P} \lor \ingamma{x}{Q} \}
  \]
\end{definition}

\begin{lemma}
  Consider any $P, Q \in \wfdomn$, and a position $i \le n$. If $\ithm{\tunion{P}{Q}} = 1$, then
  \[
  \exists x, y \in \text{concrete\_union}(P, Q) : \ith{x} \neq \ith{y}.
  \]
\end{lemma}

\begin{proof}
  The proof progresses on a case analysis of $\ithm{P}$ and $\ithm{Q}$.

  \Case{1}{$\ithm{P} = \ithm{Q} = 0$}. Since we know that
  $\ithm{\tunion{P}{Q}} = 1$, by the definition of
  \lstinline{tnum_union}, $\ithv{P} \oplus \ithv{Q} = 1$.  This means
  $\ithv{P} \neq \ithv{Q}$, and since $\ingamma{\tv{P}}{P} \land
  \ingamma{\tv{Q}}{Q}$, we pick $x = \tv{P}$ and $y = \tv{Q}$,
  satisfying our goal.

  \Case{2}{$\ithm{P} = 0, \ithm{Q} = 1$}. In this case, $\gamma({Q})$
  itself has two elements that disagree at the $i$-th bit: $x =
  \tv{Q}$ and $y = \orvm{Q}$ (membership of the latter,
  $\ingamma{y}{Q}$, is shown by
  Lemma~\ref{lem:ingamma_value_bitor_mask}).

  \Case{3}{$\ithm{P} = 1, \ithm{Q} = 0$}. Similar to the previous
  case. We pick $x = \tv{P}$ and $y = \orvm{P}$.

  \Case{4}{$\ithm{P} = \ithm{Q} = 1$}. We can reuse the choices for
  $x$ and $y$ from either Case 2 or Case 3 in this case.
\end{proof}

\section{Soundness of \lstinline{tnum_mul}}

\newcommand{\bmulshrId}{\text{bvec\_mul\_loop}}
\newcommand{\bmulshrABC}{\bmulshr{a}{b}{c}}
\newcommand{\bmulshrABCp}{\bmulshr{a'}{b'}{c'}}
\newcommand{\bmulshr}[3]{$\bmulshrId{}(#1, #2, #3)$}

\newcommand{\tmulshrId}{\text{tnum\_mul\_loop}}
\newcommand{\tmulshr}[3]{$\tmulshrId{}(#1, #2, #3)$}
\newcommand{\tmulshrABC}{\tmulshr{A}{B}{C}}
\newcommand{\tmulshrABCp}{\tmulshr{A'}{B'}{C'}}
\newcommand{\tmulId}{\text{tnum\_mul}}
\newcommand{\tmulAB}{$\tmulId{}(A, B)$}

We model the while loop from our algorithm as a recursive function
\tmulshrABC{}, where $C$ is the partial accumulator. Note that $A$ is
a well-formed tnum of length $m \ge 0$, and the other two arguments
are well-formed tnums of length $n > 0$. \tmulAB{}, where $A$ and $B$
are well-formed tnums of equal length, can be defined as
\tmulshr{A}{B}{\text{TNUM}(0, 0)}.

Apart from taking the partial accumulator as an input argument,
\tmulshrId{} employs unpadded right-shift, reducing the length of the
multiplier $A$ by one (this does not affect its value compared to the
zero-padded right shift used in the actual algorithm). It is the
truncated multiplier that gets passed to the recursive call
corresponding to the next iteration of the while loop. This truncation
and the presence of the partial accumulator as an argument helps us
present a soundness proof using induction.

Now, the main goal of the proof becomes the following:

\begin{lemma}
  Consider a well-formed tnum $A$ and a bit vector $\ingamma{a}{A}$,
  both of length $m \ge 0$. Also consider well-formed tnums $B$, $C$,
  and bit vectors $\ingamma{b}{B}, \ingamma{c}{C}$, all of length $n >
  0$. Then

  \[
  \ingamma{\bmulshrId(a, b, c)}{\tmulshrId(A, B, C)}
  \]

  where \bmulshrId{}~\footnote{Proven to be correct against the
  built-in multiplication of natural numbers in Rocq.} is a bit-vector
  analog of \tmulshrId{}.
\end{lemma}

\begin{proof}
  The proof progresses by induction on $m$, the length of $A$.

  \basecase{$m = 0.$} For a zero-length multiplier, both the tnum and
  bit vector multiplication algorithms return the partial accumulator
  as the final product. This satisfies our goal due to the
  precondition $\ingamma{c}{C}$.

  \inductionStep{} When $m > 0$, \tmulshrABC{} is \tmulshrABCp{},
  where $C'$ is the updated accumulator obtained using tnum addition
  and union as needed, and $A'$ and $B'$ are the shifted versions of
  the multiplication operands. Similarly, \bmulshrABC{} becomes
  \bmulshrABCp{}. So, we need to show the membership of \bmulshrABCp{}
  in \tmulshrABCp{}. But this will follow from the induction
  hypothesis if we can show $\ingamma{a'}{A'}$, $\ingamma{b'}{B'}$,
  and $\ingamma{c'}{C'}$. The first two memberships follow from the
  soundness of tnum shifts. But $\ingamma{c'}{C'}$ needs discussion
  because the calculation of $C$ is not straightforward.
  
  There are only two cases to consider: $\lsb{\tv{A}}
  = 1$ and $\lsb{\tm{A}} = 1$.
  In the first case, $C' = \text{tnum\_add}(C, B)$, where the bit
  vector multiplication algorithm computes $c' = \text{bvec\_add}(c,
  b)$. $\ingamma{c'}{C'}$ due to the soundness of
  \text{tnum\_add}.

  \newcommand{\newCforMu}{\text{tnum\_union}(C, \text{tnum\_add}(C, B))}
  
  When $\lsb{\tm{A}} = 1$, the tnum multiplication algorithm computes
  $C' = \newCforMu{}$. Since $\lsb{a}$ can be either $0$ or $1$ in
  this case, we need to show the following, considering two different
  possible executions of \bmulshrId{}:

  \begin{enumerate}
  \item $\ingamma{c}{\newCforMu{}}$ (case: $\lsb{a} = 0$)
  \item $\ingamma{\text{bvec\_add}(c, b)}{\newCforMu{}}$ (case:
    $\lsb{a} = 1$)
  \end{enumerate}

  Both are satisfied by the memberships given in the preconditions,
  the soundness of \text{tnum\_add}, and the soundness of
  \text{tnum\_union}.

  Since we showed that $\ingamma{a'}{A'}$, $\ingamma{b'}{B'}$, and
  $\ingamma{c'}{C'}$, our original goal follows from the induction
  hypothesis.
\end{proof}

\end{document}
